\section*{Capítulo \ref{ch:g4:what-a-fire-needs} --- A combustão: do que o fogo precisa}

\begin{solution}{exo:g4:what-a-fire-needs:1}
\omterm{def:g4:what-a-fire-needs:fuel}{Combustível}, oxigênio (do \omterm{def:g4:air-a-mixture-of-gases:air}{ar}) e calor (para começar o fogo).
\end{solution}

\begin{solution}{exo:g4:what-a-fire-needs:2}
A madeira, a cera de vela, a gasolina e o papel são \omterm{def:g4:what-a-fire-needs:fuel}{combustíveis}. A
areia, o vidro e a água não queimam.
\end{solution}

\begin{solution}{exo:g4:what-a-fire-needs:3}
O oxigênio: o cobertor impede que o \omterm{def:g4:air-a-mixture-of-gases:air}{ar} chegue até o fogo.
\end{solution}

\begin{solution}{exo:g4:what-a-fire-needs:4}
O calor: a água esfria tanto a madeira em chamas que ela não consegue
mais queimar.
\end{solution}

\begin{solution}{exo:g4:what-a-fire-needs:5}
Três destes, quaisquer: cinza, fumaça, dióxido de carbono, vapor de água.
\end{solution}

\begin{solution}{exo:g4:what-a-fire-needs:6}
O sopro traz mais \omterm{def:g4:air-a-mixture-of-gases:air}{ar}, e portanto mais oxigênio, para o \omterm{def:g4:what-a-fire-needs:fuel}{combustível}
quente, e ele volta a queimar.
\end{solution}

\begin{solution}{exo:g4:what-a-fire-needs:7}
Água: a cera que \omterm{def:g4:what-a-fire-needs:burning}{queima} fabrica vapor de água, que vira gotinhas no
vidro frio.
\end{solution}

\begin{solution}{exo:g4:what-a-fire-needs:8}
Sair, ficando abaixado por baixo da fumaça e fechando as portas depois
de passar, e então chamar os bombeiros lá de fora.
\end{solution}

\begin{solution}{exo:g4:what-a-fire-needs:9}
O lado do \omterm{def:g4:what-a-fire-needs:fuel}{combustível}: com o gás fechado, o fogo não tem mais nada para
queimar.
\end{solution}

\begin{solution}{exo:g4:what-a-fire-needs:10}
O \omterm{def:g4:what-a-fire-needs:fuel}{combustível}: sem árvores na faixa, o fogo fica sem \omterm{def:g4:what-a-fire-needs:fuel}{combustível} quando
chega até ela.
\end{solution}

\begin{solution}{exo:g4:what-a-fire-needs:11}
Não. A maior parte da madeira, junto com o oxigênio do \omterm{def:g4:air-a-mixture-of-gases:air}{ar}, virou
substâncias novas que são gases --- dióxido de carbono, vapor de água e
fumaça --- e elas foram para o \omterm{def:g4:air-a-mixture-of-gases:air}{ar}. Só a cinza fica para trás.
\end{solution}
